\documentclass[a4paper]{article}
% Español (México) con XeLaTeX: fontspec para fuentes Unicode, polyglossia
% para la división silábica y los nombres del documento en la variante mexicana.
\usepackage{fontspec}
\usepackage{polyglossia}
\setdefaultlanguage[variant=mexican]{spanish}
% Los nombres chinos van como «pinyin (original)»: xeCJK da a los caracteres
% Han su propia fuente; sin ella XeLaTeX los omite con un aviso.
% FandolSong no cubre algunos caracteres de nombres (U+73FA); WenQuanYi Zen Hei sí.
\usepackage[AutoFallBack=true]{xeCJK}
\setCJKmainfont{FandolSong-Regular.otf}
\setCJKfallbackfamilyfont{\CJKrmdefault}{WenQuanYi Zen Hei}
% polyglossia no traduce los nombres de listings.
\AtBeginDocument{\providecommand{\lstlistingname}{}\renewcommand{\lstlistingname}{Listado}%
  \providecommand{\lstlistlistingname}{}\renewcommand{\lstlistlistingname}{Índice de listados}}
\usepackage{amsmath, amssymb}
\usepackage{graphicx}
\usepackage[margin=2.5cm]{geometry}
\usepackage[most]{tcolorbox}
\usepackage{etoolbox}
\usepackage{listings}
\usepackage{hyperref}
\usepackage{booktabs}
\usepackage{subcaption}
\usepackage{float}

\newtcolorbox{knowledgebox}[1]{enhanced, colback=blue!5!white, colframe=blue!75!black, colbacktitle=blue!75!black, coltitle=white, fonttitle=\bfseries, title={#1}, attach boxed title to top left={yshift=-2mm, xshift=2mm}, boxrule=1pt, sharp corners}
\newtcolorbox{importantbox}[1]{enhanced, colback=yellow!10!white, colframe=yellow!80!black, colbacktitle=yellow!80!black, coltitle=black, fonttitle=\bfseries, title={#1}, sharp corners}
\newtcolorbox{warningbox}[1]{enhanced, colback=red!5!white, colframe=red!75!black, colbacktitle=red!75!black, coltitle=white, fonttitle=\bfseries, title={#1}, sharp corners}

\lstset{language=Python, basicstyle=\ttfamily\small, keywordstyle=\color{blue}, stringstyle=\color{red!60!black}, commentstyle=\color{green!60!black}, breaklines=true, frame=single, numbers=left, numberstyle=\tiny\color{gray}, captionpos=b, extendedchars=true}

\newcommand{\notetitle}{[LLM Agents SP25] Informal+Formal Math Reasoning — Sean Welleck}
\newcommand{\noteauthors}{Elaboradas a partir de materiales públicos del curso}
\newcommand{\notedate}{2025}
\newcommand{\videochannel}{Berkeley RDI}
\newcommand{\videopublishdate}{2025}
\newcommand{\videoduration}{aproximadamente 60 minutos}
\newcommand{\videourl}{}
\newcommand{\videocoverpath}{cover.jpg}
\newcommand{\repourl}{https://github.com/hqhq1025/ai-course-notes}


% Footer with repo info
\usepackage{fancyhdr}
\pagestyle{fancy}
\fancyhf{}
\fancyfoot[C]{\thepage}
\fancyfoot[R]{\footnotesize\color{gray}\href{\repourl}{AI Course Notes}}
\renewcommand{\headrulewidth}{0pt}
\renewcommand{\footrulewidth}{0pt}

\begin{document}
\begin{titlepage}
\centering
{\Large Notas del curso\par}\vspace{1.2cm}
{\huge\bfseries \notetitle\par}\vspace{0.8cm}
{\large \noteauthors\par}\vspace{0.3cm}
{\large \notedate\par}\vspace{1.2cm}
\ifdefempty{\videocoverpath}{}{\includegraphics[width=0.82\textwidth,height=0.45\textheight,keepaspectratio]{\videocoverpath}\par}
\vfill
\begin{tcolorbox}[width=0.9\textwidth, colback=black!2!white, colframe=black!60, sharp corners]
\textbf{Autor o canal del video}: \videochannel\par
\textbf{Fecha de publicación}: \videopublishdate\par
\textbf{Duración del video}: \videoduration\par
\ifdefempty{\videourl}{}{\textbf{Enlace del video}: \href{\videourl}{\nolinkurl{\videourl}}\par}\par
\textbf{Página del proyecto}: \href{\repourl}{\nolinkurl{\repourl}}\par
\end{tcolorbox}
\end{titlepage}
\tableofcontents
\newpage

\section{Introducción: el puente entre el razonamiento matemático informal y el formal}

Esta clase la imparte Sean Welleck, de CMU, quien examina cómo aprovechar la IA para cerrar la brecha entre las \textbf{matemáticas informales} (informal math) y las \textbf{matemáticas formales} (formal math), combinando las ventajas de ambas.

\begin{importantbox}{Comparación entre las dos representaciones matemáticas}
\textbf{Matemáticas informales}: se expresan en lenguaje natural (texto, LaTeX, imágenes, etc.), son sumamente flexibles, pero no garantizan la corrección. El LLM puede cometer errores sutiles en los pasos intermedios.\\
\textbf{Matemáticas formales}: se escriben en un lenguaje de programación (Lean, Coq, Isabelle); que el código compile equivale a que la demostración sea correcta. Ofrecen una garantía absoluta de corrección, pero el umbral para escribirlas es alto y el proceso es laborioso.
\end{importantbox}

\subsection{Introducción a los asistentes de demostración formal (Lean)}

Lean es un asistente de demostración interactivo (Interactive Theorem Prover); sus conceptos centrales son:
\begin{itemize}
    \item \textbf{Enunciado del teorema}: expresar en lenguaje formal la proposición que se debe demostrar
    \item \textbf{Estado de la demostración} (Proof State): el contexto actual (las hipótesis) y el objetivo pendiente
    \item \textbf{Táctica} (Tactic): la regla de inferencia que se aplica en cada paso, la cual modifica el estado de la demostración
    \item \textbf{QED}: la demostración se completa cuando se resuelven todos los objetivos
\end{itemize}

\begin{knowledgebox}{Mathlib: la infraestructura de las matemáticas formales}
Mathlib es el proyecto de biblioteca matemática de Lean, construido en conjunto por numerosos colaboradores con las definiciones y los teoremas fundamentales de las matemáticas. Permite que las demostraciones formales reutilicen resultados ya existentes; el matemático Terence Tao ya ha logrado formalizar en Lean los teoremas de sus artículos.
\end{knowledgebox}

\subsection{La importancia de las matemáticas formales para la IA}

\begin{enumerate}
    \item \textbf{Salida verificable}: evita el razonamiento matemático incorrecto
    \item \textbf{Señal de recompensa perfecta}: se puede usar para el entrenamiento con RL, ya que si la demostración compila, es correcta
    \item \textbf{Banco de pruebas para el razonamiento}: abarca problemas de razonamiento desde sencillos hasta arbitrariamente difíciles
\end{enumerate}

La demostración formal de teoremas ha avanzado rápidamente en los últimos años: desde GPT-f (2020) hasta DeepSeek-Prover, los resultados en el banco de pruebas mini-F2F han subido de manera sostenida.
\section{Primera parte: incorporar el pensamiento informal a la demostración formal — LeanSTaR}

\subsection{Motivación: más allá del entrenamiento puramente formal}

Los métodos tradicionales entrenan el modelo solo con datos de demostración formal (estado $\to$ política del siguiente paso), pero esto puede no bastar para aprender el \textbf{proceso de pensamiento} necesario para producir el código.

\begin{importantbox}{La idea central de LeanSTaR}
Entrenar al modelo para que, antes de cada paso de demostración formal, primero genere un \textbf{pensamiento informal} (informal thought) de forma libre:
\begin{itemize}
    \item El modelo puede descomponer el problema, planear el siguiente paso y ejecutar cálculos informales
    \item El pensamiento ayuda a diversificar el espacio de búsqueda
    \item De forma similar a Chain-of-Thought, aumenta la capacidad expresiva del modelo
\end{itemize}
\end{importantbox}

\subsection{El algoritmo de LeanSTaR}

Entrenamiento en dos etapas:

\textbf{Etapa uno: inicialización}
\begin{enumerate}
    \item A partir de demostraciones existentes (como Mathlib), dado el estado y la política del siguiente paso, se le pide al LLM que genere retrospectivamente el pensamiento que describe esa transición
    \item Se recopilan cientos de miles de muestras (estado, pensamiento, política)
    \item Se entrena un modelo inicial para que aprenda a generar el pensamiento antes de la política
\end{enumerate}

\textbf{Etapa dos: aprendizaje por refuerzo (Expert Iteration)}
\begin{enumerate}
    \item El modelo genera en paralelo múltiples demostraciones completas
    \item Lean verifica si la demostración es correcta
    \item Las trayectorias exitosas se recopilan y se agregan al conjunto de entrenamiento
    \item Se entrena un nuevo modelo sobre el conjunto de datos ampliado
    \item Se itera este proceso
\end{enumerate}

\begin{warningbox}{La dificultad de la búsqueda en árbol al incorporar el pensamiento}
La Best-First Search tradicional funciona mal una vez que se incorpora el pensamiento informal: es difícil puntuar el pensamiento con precisión. LeanSTaR opta en cambio por una estrategia sencilla de \textbf{muestreo en paralelo}: se generan varias rutas en paralelo, se reintenta cuando hay un error y se detiene en cuanto una tiene éxito.
\end{warningbox}

\subsection{Resultados experimentales}

\begin{itemize}
    \item En el punto de referencia mini-F2F, LeanSTaR llevó al modelo de estar por debajo de la referencia del agente GPT-4 a superarla
    \item \textbf{Hallazgo clave}: el modelo con pensamiento se beneficia más de la escalabilidad del presupuesto de inferencia (scaling inference budget), porque el pensamiento diversifica de manera efectiva el espacio de búsqueda
    \item En el pensamiento informal, el modelo ejecutó pasos de razonamiento que no existen en el código formal
\end{itemize}

Desde entonces, el pensamiento informal combinado con entrenamiento por RL ha sido adoptado ampliamente por la comunidad de demostración de teoremas (DeepSeek-Prover, entre otros), y también es el paradigma central de modelos de razonamiento como O1 y DeepSeek-R1.

\subsection{Resumen de la sección}

La idea central de LeanSTaR: entrenar solo con código formal puede no bastar para aprender el proceso de pensamiento necesario para generar código. Aprender a generar pensamiento informal mediante RL puede mejorar de manera significativa la capacidad de demostración.
\section{Segunda parte: demostración informal — borrador-esbozo-demostración}

\subsection{Motivación: combinar el razonamiento de alto nivel con la demostración de bajo nivel}

\begin{knowledgebox}{Sledgehammer: herramienta de demostración automática}
Sledgehammer en Isabelle es una herramienta de «martillo» (Hammer) que envía el problema a demostradores de teoremas automáticos externos (lógica de primer orden, lógica de orden superior, resolutores SMT); es muy hábil para llenar pequeños vacíos en una demostración, pero no puede manejar estrategias de demostración de alto nivel que requieren creatividad.
\end{knowledgebox}

\subsection{El método Draft, Sketch, Prove (DSP)}

\begin{importantbox}{Proceso de demostración en tres pasos}
\begin{enumerate}
    \item \textbf{Draft} (borrador): el LLM genera una demostración informal (en lenguaje natural)
    \item \textbf{Sketch} (esbozo): el LLM traduce la demostración informal a un esbozo formal, dejando vacíos por llenar
    \item \textbf{Prove} (demostración): un demostrador de bajo nivel (como Sledgehammer) llena todos los vacíos
\end{enumerate}
Si se llenan todos los vacíos, se obtiene una demostración formal completa.
\end{importantbox}

La estrategia de búsqueda difiere de la búsqueda en árbol tradicional: se generan múltiples borradores informales o múltiples traducciones a esbozos formales, se intenta llenar los vacíos, y basta con que uno tenga éxito para obtener la demostración.

\subsection{Resultados experimentales}

En mini-F2F (en la época de CodeX/Minerva):
\begin{itemize}
    \item El desempeño mejora notablemente conforme aumenta el número de muestreos
    \item Con muestreos suficientes, las demostraciones informales generadas por el LLM incluso superan el desempeño obtenido con demostraciones escritas a mano por humanos
\end{itemize}

\section{Tercera parte: construir un Hammer para Lean --- LeanHammer}

\subsection{Arquitectura del pipeline de Hammer}

\begin{importantbox}{Componentes de LeanHammer}
\begin{enumerate}
    \item \textbf{Selección de premisas} (Premise Selection): un modelo neuronal de recuperación aplica embedding al objetivo y a las premisas candidatas, y usa similitud coseno para seleccionar los teoremas, definiciones o lemas más relevantes
    \item \textbf{Capa de traducción} (lean-auto): traduce la teoría de tipos dependientes de Lean a la lógica del demostrador automático
    \item \textbf{Demostrador automático de teoremas} (Zipperposition, entre otros)
    \item \textbf{Reconstrucción de la demostración} (Duper): traduce la salida del demostrador automático de vuelta a una demostración en Lean
    \item \textbf{Búsqueda en árbol} (AESOP): organiza todo el proceso de búsqueda
\end{enumerate}
\end{importantbox}

El usuario solo necesita invocar el comando \texttt{hammer} dentro de la demostración, y el sistema realiza automáticamente la selección de premisas, la traducción, la demostración automática y la reconstrucción.

\subsection{Resultados experimentales}

\begin{itemize}
    \item El selector neuronal de premisas de LeanHammer (con solo 100 millones de parámetros) supera claramente a los trabajos anteriores
    \item Es complementario al selector simbólico de premisas (Moogle); al tomar la unión de ambos, el desempeño mejora aún más
    \item El desempeño se está acercando al límite teórico superior que se obtiene al usar las premisas reales presentes en demostraciones humanas
\end{itemize}

\subsection{Resumen de la sección}

DSP y LeanHammer representan dos enfoques complementarios: el primero guía la estructura de la demostración mediante bosquejos de alto nivel, mientras que el segundo ofrece una herramienta de automatización de bajo nivel, poderosa, para completar los vacíos.
\section{Cuarta parte: hacia la asistencia de IA para matemáticas de nivel de investigación}

\subsection{Los retos particulares de la formalización de nivel de investigación}

La formalización de matemáticas de nivel de investigación (como el proyecto PFR de Terence Tao) implica:
\begin{enumerate}
    \item Crear un \textbf{plano} (Blueprint): descomponer la demostración en un grafo de dependencias
    \item Escribir una gran cantidad de \textbf{definiciones nuevas} y \textbf{lemas auxiliares}
    \item Ensamblar el teorema final y su demostración
\end{enumerate}

\begin{warningbox}{Limitaciones de las pruebas de referencia actuales}
Las pruebas de referencia más usadas (como mini-F2F) se basan en problemas de competencias matemáticas, que se caracterizan por ser autocontenidos y depender de resultados estándar. En cambio, los teoremas de los proyectos de investigación reales dependen de \textbf{definiciones y lemas nuevos propios del proyecto}; esta brecha hace que los modelos con buen desempeño en problemas de competencia puedan fallar gravemente en proyectos reales.
\end{warningbox}

\subsection{La prueba de referencia MiniCTX}

\begin{importantbox}{MiniCTX: una prueba de referencia de demostración de teoremas dependiente del contexto}
\begin{itemize}
    \item Extrae teoremas de prueba de proyectos reales de Lean (como Mathlib y PFR)
    \item Selecciona los teoremas por una fecha de corte, y puede actualizarse continuamente para mantenerse por delante de la fecha de corte del entrenamiento de los modelos
    \item Distingue entre dependencias dentro del mismo archivo y dependencias entre archivos
    \item Los experimentos muestran que: en problemas de competencia, el contexto casi no afecta el desempeño; en proyectos reales, \textbf{la falta de contexto provoca una caída abrupta del desempeño}
\end{itemize}
\end{importantbox}

\subsection{Herramienta práctica: LLMLean}

LLMLean es un complemento de Lean que puede conectarse a modelos locales (mediante Ollama o LM Studio) o a modelos de API (OpenAI, Claude, Gemini, etc.) y ofrece directamente en VS Code sugerencias de demostración verificadas por Lean.
\section{Síntesis y ampliación}

\subsection{Tres caminos de conexión}

\begin{enumerate}
    \item \textbf{Razonamiento informal + demostración formal} (LeanSTaR): entrena al modelo para razonar de manera informal antes de cada paso de la demostración
    \item \textbf{Demostración informal + relleno formal} (DSP + LeanHammer): primero genera un bosquejo de alto nivel y luego usa herramientas automatizadas para rellenar los detalles
    \item \textbf{Asistencia a nivel de investigación} (MiniCTX + LLMLean): demostración consciente del contexto y herramientas prácticas orientadas a proyectos matemáticos reales
\end{enumerate}

\subsection{Lecturas adicionales}
\begin{itemize}
    \item LeanSTaR (ICLR 2025): entrenamiento con RL de un modelo que genera razonamiento informal
    \item Draft, Sketch, Prove (ICLR 2023): método de demostración por bosquejo y relleno
    \item MiniCTX (ICLR 2025 Oral): referencia de demostración de teoremas dependiente del contexto
    \item LLMLean: \url{https://github.com/cmu-l3/llmlean}
    \item DeepSeek-Prover: demostrador de vanguardia que integra el razonamiento informal
\end{itemize}

\end{document}
